% !TEX root = ../linnik399.tex
\section{Prime detection}\label{sec:detection}

Our goal is to make a nonnegative weighted sum over primes in a reduced residue class
strictly positive. The explicit formula expresses this as a comparison between a main
term and a sum over zeros. We first fix the weight, then state the precise zero-sum
bound that will suffice.

\subsection{The kernel}\label{subsec:kernel}
Fix
\[
  T=0.416829,\qquad \kappa=T/16,
\]
and the sixteen step heights $b_0,\dots,b_{15}$ of \cref{tab:beta}. Put
\[
  \psi(t)=\sum_{i=0}^{15}b_i\,\1_{[i\kappa,(i+1)\kappa)}(t),\qquad
  \Psi(z)=\int_0^T\psi(t)e^{-zt}\,dt .
\]
For the target exponent $L$ put $A=L-2T$; at $L=3.99$ we have $A=3.156342$. The \emph{prime
weight} is
\[
  h(t)=(\psi*\psi)(t-A),
\]
which is supported in $[A,L]$. Its Laplace transform is
\begin{equation}\label{eq:Hkernel}
  H(z)=\int h(t)e^{-zt}\,dt=e^{-Az}\,\Psi(z)^2,\qquad H_0=H(0)=\Psi(0)^2 .
\end{equation}
Numerically $\Psi(0)=\kappa\sum_ib_i=0.193642\ldots$ and $H_0=0.037497\ldots$.

\begin{table}[ht]
\centering
\small
\begin{tabular}{rl@{\qquad}rl@{\qquad}rl@{\qquad}rl}
\toprule
$i$ & $b_i$ & $i$ & $b_i$ & $i$ & $b_i$ & $i$ & $b_i$ \\
\midrule
0 & 0.02318741 & 4 & 0.23237155 & 8 & 0.47877833 & 12 & 0.75172165 \\
1 & 0.07159337 & 5 & 0.29081583 & 9 & 0.54505502 & 13 & 0.82236428 \\
2 & 0.12265085 & 6 & 0.35147264 & 10 & 0.61279917 & 14 & 0.89340300 \\
3 & 0.17627704 & 7 & 0.41418551 & 11 & 0.68177389 & 15 & 0.96451862 \\
\bottomrule
\end{tabular}
\caption{The step heights of the kernel $\psi$. They are exact decimal numbers.}
\label{tab:beta}
\end{table}

Xylouris \cite[(3.51)]{Xylouris2011nullstellen}, following Heath-Brown \cite[\S13]{HeathBrown1992zero},
uses the triangular weights
\[
  h_{L',K}(t)=\begin{cases}
    t-(L'-2K), & L'-2K\le t\le L'-K,\\
    L'-t, & L'-K\le t\le L',\\
    0, & \text{otherwise},
  \end{cases}
\]
for $L'>2K>0$; the transform of $h_{L',K}$ is $e^{-(L'-2K)z}\bigl((1-e^{-Kz})/z\bigr)^2$.

\begin{lemma}[Decomposition into positive triangular weights]\label{lem:triangles}
For the step function $\psi$ and shifted convolution $h$ defined above,
\[
  h=\sum_{k=0}^{30}c_k\,h_{A+(k+2)\kappa,\,\kappa},\qquad
  c_k=\sum_{\substack{0\le i,j\le15\\ i+j=k}}b_ib_j>0 .
\]
Every triangle in this sum has left endpoint $A+k\kappa\ge A$.
\end{lemma}

\begin{proof}
For $i,j\ge0$ the convolution $\1_{[i\kappa,(i+1)\kappa)}*\1_{[j\kappa,(j+1)\kappa)}$ is the tent
function of height $\kappa$ supported on $[(i+j)\kappa,(i+j+2)\kappa]$, which is
$h_{(i+j+2)\kappa,\kappa}$. Expanding $\psi*\psi$ and shifting by $A$ gives the formula. Each
$c_k$ is a sum of products of positive numbers, and the sum over $i+j=k$ is nonempty for
$0\le k\le30$.
\end{proof}

\subsection{The positivity criterion}
The following is Xylouris's form of the explicit formula for a finite combination of triangles
\cite[(3.56)--(3.58) and the lines between (3.56) and (3.57), pp.~34--35]{Xylouris2011nullstellen}; it
refines Heath-Brown's Lemmas~13.1 and~13.2 \cite{HeathBrown1992zero}.

\begin{inp}[Xylouris's criterion]\label{inp:criterion}
Let $h=\sum_i\alpha_ih_{L_i,K_i}$ be a finite real combination of triangles with
$L_i>2K_i+3$ for every $i$, and let $H$ be its Laplace transform. For every $\eps>0$ there are
$C_0=C_0(\eps)$ and $q_0=q_0(\eps)$ such that for every $q\ge q_0$ and every $a$ coprime to $q$,
\[
  \sum_{p\equiv a\,(q)}\frac{\log p}{p}\,h\Bigl(\frac{\log p}{\LL}\Bigr)\ \ge\
  \frac{\LL}{\varphi(q)}\Bigl[H(0)-\sum_{\chi\ne\chi_0}\ {\sum_\rho}'\ \bigl|H\bigl((1-\rho)\LL\bigr)\bigr|
  -\eps\Bigr],
\]
where $\sum'_\rho$ runs over the zeros $\rho=\beta+i\gamma$ of $L(s,\chi)$, with multiplicity, in the
square $1-C_0/\LL\le\beta\le1$, $|\gamma|\le C_0/\LL$.
\end{inp}

Xylouris states the extension to combinations without a separate proof, and it follows from
Heath-Brown's proof of his Lemma~13.2 \cite[\S13]{HeathBrown1992zero}. That proof starts from the explicit
formula \cite[Lemma~13.1]{HeathBrown1992zero}, which is linear in the weight, and bounds the zeros
outside the square by summing over rectangles $1-\frac{m+1}\LL\le\beta\le1-\frac m\LL$,
$\frac n{2\LL}\le|\gamma|\le\frac n\LL$ with $\max(m,n)\ge R$, using a zero-density bound $\ll e^{3m}(1+n/\LL)^{3/2}$
for each rectangle and $|F((1-\rho)\LL)|\ll e^{-(L'-2K)m}n^{-2}$ for each zero. For a combination with positive
coefficients $c_k$ one has $|H|\le\sum_kc_k|H_k|$, so the same estimate holds with $L'-2K$ replaced by the
smallest left endpoint, which exceeds $3$; and the zeros in the square contribute at most
$\sum|H((1-\rho)\LL)|$.

In the normalized coordinates \eqref{eq:normalized} the square is
$\{0\le\lambda\le C_0,\ |\mu|\le C_0\}$. Enlarging it only adds nonnegative terms to the zero
sum, so the conclusion holds for every square $\{\lambda\le C_P,\,|\mu|\le C_P\}$ with $C_P\ge
C_0$. We fix
\begin{equation}\label{eq:CP}
  C_P\ge\max\bigl(C_0(\eta H_0),3\bigr),
\end{equation}
which fixes the prime window $R_P=\{0\le\lambda\le C_P,\ |\mu|\le C_P\}$ of \cref{sec:notation}. Since $C_P$ is fixed and
$\log\log\LL\to\infty$, for large $q$ every zero in $R_P$ lies in the rectangle $R(l)$ of
\cref{sec:notation}.

\begin{definition}[Normalized zero sum]\label{def:zerosum}
For the kernel $H$ in \eqref{eq:Hkernel} and the prime window fixed by \eqref{eq:CP}, put
\[
  W=W(q)=\frac1{H_0}\sum_{\chi\ne\chi_0}\ \sum_{\rho\in R_P}\bigl|H\bigl((1-\rho)\LL\bigr)\bigr| ,
\]
where each zero is counted with its multiplicity.
\end{definition}

\begin{proposition}[A zero-sum bound that guarantees a prime]\label{prop:detect}
Let $L>3+2T$. There is $q_0$ such that for every $q\ge q_0$ with $W(q)\le1-2\eta$ and every $a$ coprime
to $q$ there is a prime $p$ with
\[
  p\equiv a\pmod q,\qquad q^{A}<p<q^{L},\qquad A=L-2T.
\]
Here $W(q)$ is formed from the kernel and prime window for this fixed $L$; the threshold
$q_0$ is independent of $a$ and of the zero configuration of $q$.
\end{proposition}

\begin{proof}
By \cref{lem:triangles}, $h$ is a finite positive combination of triangles $h_{L_i,\kappa}$ with
$L_i-2\kappa=A+k\kappa\ge A>3$. So \cref{inp:criterion} applies with $\eps=\eta H_0$; it
gives a $q_0$ such that for every $q\ge q_0$ with $W(q)\le1-2\eta$,
\[
  \sum_{p\equiv a\,(q)}\frac{\log p}{p}\,h\Bigl(\frac{\log p}{\LL}\Bigr)\ge
  \frac{\LL H_0}{\varphi(q)}\bigl(1-W-\eta\bigr)\ge\frac{\LL H_0\,\eta}{\varphi(q)}>0 .
\]
The sum is therefore nonempty. Since $h$ vanishes outside $(A,L)$, some prime $p\equiv a\pmod q$
satisfies $A<\log p/\LL<L$.
\end{proof}

At $L=3.99$ the condition $L>3+2T=3.833658$ holds. Two features of \cref{inp:criterion} matter
below.
\begin{enumerate}[(i)]
  \item \emph{Occurrence form.} The zero sum charges each zero occurrence separately, through the
    absolute value of its own term. Every bound in \cref{sec:costs} is a bound for a sum of such
    individual terms. In particular we never need the cancellation between zeros that a grouped
    form of the explicit formula would record.
  \item \emph{Size of a term.} Since $|\Psi(\lambda-i\mu)|\le\Psi(\lambda)$ for $\lambda\ge0$,
    \begin{equation}\label{eq:termbound}
      \bigl|H\bigl((1-\rho)\LL\bigr)\bigr|=e^{-A\lambda}\,|\Psi(\lambda-i\mu)|^2\le
      e^{-A\lambda}\,\Psi(\lambda)^2\le H_0e^{-A\lambda}.
    \end{equation}
    A zero at distance $\lambda$ from $\sigma=1$ costs at most $e^{-A\lambda}$ times $H_0$. The
    difficulty is entirely with the number of zeros close to $\sigma=1$.
\end{enumerate}

If no nonprincipal $L(s,\chi)$ has a zero in $R(l)$, then for large $q$ it has none in $R_P$, and
$W=0$. This is the first exterior regime (\cref{sec:exteriors}).
